% SGA 3, Expose V, Loop 1 English reconstruction.
% Scope: section 4 heading, Theorem 4.1, editor's notes 16--20, and the
% paragraph announcing the proof strategy.
% Authority: Polo--Gille Exp5-13oct24.pdf, component-local p. 11
% (combined-reader p. 273). This re-edited page bears printed header 261
% and internal original-page markers 263--264.
% Jacob Reinhold's CC BY 4.0 Markdown is comparison/drafting material only.

\section*{4. 对有限平坦群胚取商 (一个特殊情形的证明)}
\phantomsection\label{sec:sga3-V-4}
\addcontentsline{toc}{section}{4. 对有限平坦群胚取商}

\medskip
\phantomsection\label{thm:sga3-V-4.1}
\noindent\textbf{定理 4.1.}---
\textit{考虑一个 \((\mathbf{Sch}/S)\)-群胚}

% Native Loop-2 reconstruction; local p. 11 / combined p. 273.
\begin{equation*}
  \resizebox{0.40\textwidth}{!}{%
    \SGAThreeNativeInput{figures/exp5/p261_theorem41_groupoid.png}}
  \phantomsection\label{fig:sga3-exp5-theorem41-groupoid}
\end{equation*}

\textit{并假设以下条件成立:}%
\footnote{\label{note:sga3-V-16}编者注: 由于 \(d_0=d_1\sigma\),
其中 \(\sigma\) 是 \(X_1\) 的对合自同构,
这两个条件关于 \(d_1\) 和 \(d_0\) 对称.
此外, \(d_0d_1^{-1}(x)=d_1d_0^{-1}(x)\).}
\begin{enumerate}
\item[\textit{a)}]\phantomsection\label{thm:sga3-V-4.1-a}%
  \textit{\(d_1\) 有限且局部自由;}
\item[\textit{b)}]\phantomsection\label{thm:sga3-V-4.1-b}%
  \textit{对每个 \(x\in X_0\), 集合 \(d_0d_1^{-1}(x)\)
  包含在 \(X_0\) 的某个仿射开子集中.}%
  \footnote{\label{note:sga3-V-17}编者注:
  \hyperref[thm:sga3-V-4.1-b]{条件~b)} 不能省去.
  H.~Hironaka 给出了有限群 \(G=\mathbb Z/2\mathbb Z\)
  在一个复完备簇 \(X_0\) 上作用的例子, 其商 \(X_0/G\)
  是一个不是概形的代数空间; 参见
  \(\hyperref[bib:sga3-V-Hi62]{[{\rm Hi62}]}\), 以及
  \(\hyperref[bib:sga3-V-Mum65]{[{\rm Mum65}]}\), 第~4~章, \S3.}
\end{enumerate}
\textit{则:}
\begin{enumerate}
\item[(i)]\phantomsection\label{thm:sga3-V-4.1-i}%
  \textit{在 \((\mathbf{Sch}/S)\) 中存在 \((d_0,d_1)\)
  的余核 \((Y,p)\). 此外, 这样的二元组 \((Y,p)\)
  也是所有赋环空间构成的范畴中 \((d_0,d_1)\) 的余核;}
\item[\SGAInlineRefTarget{thm:sga3-V-4.1:item:ii}(ii)]\phantomsection\label{thm:sga3-V-4.1-ii}%
  \textit{\(p\) 是整的开态射, 且当 \(X_0\) 仿射时, \(Y\) 也仿射;}%
  \footnote{\label{note:sga3-V-18}编者注: \(p\) 为开态射的断言是新增的,
  其论证与 \hyperref[lem:sga3-V-6.1]{6.1} 中所给的类似论证相同.}
\item[(iii)]\phantomsection\label{thm:sga3-V-4.1-iii}%
  \textit{分量为 \(d_0\) 和 \(d_1\) 的态射
  \[
    X_1\longrightarrow X_0\times_YX_0
  \]
  是满射;}
\item[(iv)]\phantomsection\label{thm:sga3-V-4.1-iv}%
  \textit{若 \((d_0,d_1)\) 是等价对, 则
  \(X_1\to X_0\times_YX_0\) 是同构}%
  \footnote{\label{note:sga3-V-19}编者注: 此时
  \(X_1\to X_0\times_SX_0\) 因而是浸入
  (EGA~I, 5.3.10); 另见第~VI\textsubscript{B}~篇, 9.2.1.
  在较弱的假设, 即 \(d_0\) 和 \(d_1\) 有限而未必平坦的条件下,
  关于概形范畴或代数空间范畴中商的存在性, 参见
  \(\hyperref[bib:sga3-V-An73]{[{\rm An73}]}\), \S1.1,
  \(\hyperref[bib:sga3-V-Fe03]{[{\rm Fe03}]}\),
  \(\hyperref[bib:sga3-V-Ko08]{[{\rm Ko08}]}\) 及其中的参考文献.}
  \textit{并且 \(p:X_0\to Y\) 有限且局部自由.}%
  \footnote{\label{note:sga3-V-20}编者注: 下述推论已明确写出;
  参见 \(\hyperref[bib:sga3-V-Ray67a]{[{\rm Ray67a}]}\),
  定理~1(iii), 以及 \hyperref[sec:sga3-V-5]{第~5~节} 末尾
  \hyperref[proof:sga3-V-5-b-general]{给出的证明}.}
  \textit{此外, \((Y,p)\) 是 fppf 拓扑的层范畴中
  \((d_0,d_1)\) 的余核, 且对每个基变换 \(Y'\to Y\),
  概形 \(Y'\) 是从 \(X_\ast\) 沿基变换
  \(X_0\times_YY'\to X_0\) 得到的群胚
  \(X_\ast\times_YY'\) 的余核.}
\end{enumerate}
\textit{特别地, 对每个基变换 \(S'\to S\),
概形 \(Y'=Y\times_SS'\) 是 \(S'\)-群胚
\(X'_\ast=X_\ast\times_SS'\) 的余核.
因此, 在这种情形下, 商的构造与基变换可交换.}

由 \hyperref[thm:sga3-V-4.1-i]{(i)} 立即可知,
\(Y\) 的底拓扑空间, 是 \(X_0\) 的底拓扑空间
对 \((\mathbf{Sch}/S)\)-群胚 \(X_\ast\) 所定义的等价关系的商.

先在 \(X_0\) 仿射且 \(d_1\) 局部自由、秩为常数 \(n\) 的情形下
证明 \hyperref[thm:sga3-V-4.1]{本定理}.
然后再说明如何将一般情形化归到这一特殊情形.
% SGA3-V-V2-EOL-PRESERVATION
% Translation statistics: source_lines=98; output_lines=96; inline_math=55; display_math=1; equation=0; equation_star=1; align=0; footnotes=5.
% Review: complete assigned component; pending_TODO=0; all math symbols, labels, references, diagram paths and environments checked; formula prose translated.
